\documentclass[14pt]{extarticle}
\usepackage{preamble}
\geometry{letterpaper, landscape, margin=0.5in}
\usepackage{qrcode}
\renewcommand{\answer}[1]{}
\setlength{\parskip}{4pt}

\begin{document}
\pagestyle{empty}

\daybanner{Unit 2 \textperiodcentered\ Topic 2.11 \textperiodcentered\ TODAY'S PROBLEM}{Comparing a Score of 92}

\noindent\begin{minipage}[t]{0.57\linewidth}
\vspace{0pt}
Suppose exam scores at School A and School B are each approximately normal. Both distributions have mean $80$ points, but School A has standard deviation $9$ points and School B has standard deviation $15$ points.\par\smallskip A counselor says, ``A score of 92 is at about the 90th percentile at both schools because both schools have the same mean.'' A data coach says, ``The different spreads could put 92 at different percentiles.''\par
\end{minipage}\hfill
\begin{minipage}[t]{0.40\linewidth}
\vspace{0pt}
\begin{center}
\begin{tikzpicture}
\begin{axis}[
  width=0.88\linewidth,
  height=1.75in,
  xmin=35, xmax=125,
  ymin=0, ymax=0.050,
  axis x line=bottom,
  axis y line=left,
  ytick=\empty,
  xtick={50,65,80,92,110},
  xlabel={Exam score (points)},
  tick label style={font=\small},
  label style={font=\small},
  legend style={font=\small,draw=none,fill=none,at={(0.5,1.02)},anchor=south}
]
\addplot[dayblue,very thick,domain=35:125,samples=181]
  {1/(9*sqrt(2*pi))*exp(-0.5*((x-80)/9)^2)};
\addplot[warmred,very thick,dashed,domain=35:125,samples=181]
  {1/(15*sqrt(2*pi))*exp(-0.5*((x-80)/15)^2)};
\addplot[black,densely dotted] coordinates {(92,0) (92,0.049)};
\legend{{School A: $\mu=80,\ \sigma=9$},{School B: $\mu=80,\ \sigma=15$}}
\end{axis}
\end{tikzpicture}
\end{center}
\end{minipage}

\begin{firsttakebox}{FIRST TAKE (5 min, on your sheet)}
Would a 92 stand out equally at both schools? Give one reason from the picture.
\end{firsttakebox}

\begin{description}[leftmargin=1.15in, labelwidth=1in, itemsep=2pt, topsep=2pt]
  \item[\dokbadge{1}~(a)] Identify which school has the more variable score distribution and state how the corresponding normal curve shows that difference.
  \item[\dokbadge{2}~(b)] For a score of 92, calculate the $z$-score and percentile at each school. Show the standardization and report each percentile to the nearest tenth of a percent.
  \item[\dokbadge{3}~(c)] In 3--4 sentences, decide whose claim is better supported. Use both percentiles from (b) and explain how the different spreads change what a score of 92 means at each school.
\end{description}

\vfill
\noindent\begin{minipage}{0.84\linewidth}
\textbf{Video + follow-along:} \href{https://robjohncolson.github.io/apstats-live-worksheet/u5_lesson1-2_live.html}{\texttt{\detokenize{u5_lesson1-2_live.html}}} (scan the code)\par
\textbf{Finish (a)--(c) after the video. Turn in your sheet.}
\end{minipage}\hfill
\qrcode[height=0.75in]{https://robjohncolson.github.io/apstats-live-worksheet/u5_lesson1-2_live.html}

\end{document}
